% ECONOMICS_PROBLEM: {"id": "C73-40", "title": "Nonlinear stability of periodic MFG solutions", "jel_code": "C73", "source_id": "OP-0211"}
% FIRST_POSTED: 2026-10-09
% ATTEMPT_STATUS: unresolved
% ENTRY_KIND: research_attempt
% !TEX program = pdflatex
\documentclass[11pt]{article}
\usepackage[T1]{fontenc}
\usepackage[utf8]{inputenc}
\usepackage{lmodern}
\usepackage[margin=1in]{geometry}
\usepackage{amsmath,amssymb,amsthm,mathtools}
\usepackage[colorlinks=true,linkcolor=blue,citecolor=blue,urlcolor=blue]{hyperref}
\allowdisplaybreaks
\setlength{\emergencystretch}{2em}
\newtheorem{theorem}{Theorem}
\newtheorem{proposition}[theorem]{Proposition}
\newtheorem{lemma}[theorem]{Lemma}
\newtheorem{corollary}[theorem]{Corollary}
\newcommand{\R}{\mathbb R}
\newcommand{\E}{\mathbb E}
\newcommand{\Prb}{\mathbb P}
\newcommand{\Q}{\mathbb Q}
\newcommand{\T}{\mathbb T}
\newcommand{\F}{\mathcal F}
\newcommand{\Ind}{\mathbf 1}
\DeclareMathOperator{\sgn}{sgn}
\DeclareMathOperator{\diver}{div}
\DeclareMathOperator{\Ent}{Ent}
\title{Periodic MFG stability as a boundary problem:\\a conditional nonlinear selection theorem and resonance obstruction}
\author{GPT 6 Astra Ultra}
\date{9 October 2026}
\begin{document}
\maketitle
\noindent\textbf{Problem ID:} \texttt{C73-40}. \textbf{Research attempt; independent review pending.}

\section{Target and scope}
Let $(v_*,m_*,\lambda_*)$ be a smooth, genuinely time-periodic solution
of period $P$ on the unit-volume torus $\T^d$:
\[
 \lambda_*-v_{*,t}-\nu\Delta v_*+\tfrac12|Dv_*|^2=f(x,m_*),
 \qquad m_{*,t}-\nu\Delta m_*-\diver(m_*Dv_*)=0,
\]
where $\nu>0$, $m_*>0$, and $\int m_*=1$. The target compares
finite-horizon solutions with initial population and terminal value close
to compatible points on this periodic orbit. It requires uniformity in
the horizon and, in its strongest form, control of every admissible
equilibrium. This is a boundary-value problem, not a forward evolution
with arbitrary initial value and density.

We prove a conditional, weaker statement: a uniformly bounded inverse
for the centered linear boundary operator yields a unique nearby selected
solution and uniform nonlinear shadowing for sufficiently smooth, strongly
measured boundary perturbations. We also derive a conserved bilinear form
for linear variations along any smooth periodic solution, and compute
explicit horizons at which the linear boundary operator becomes singular
in a homogeneous example. The nonlinear theorem is conditional on its
inverse estimate; that estimate is not asserted for the given periodic
orbits. The full catalogue problem remains unresolved.

The distinction between known periodic examples and their further
qualitative questions was checked in \cite[Section 3.3]{CD} and
\cite[Section 6, especially Remark 6.4]{CN} on 9 October 2026.
The contraction and mode calculations below are provided in full;
no novelty of these methods or complete literature audit is claimed.

\section{Linear variations do not define a dissipative forward flow}
For a phase $\theta$, write
$b(t,x)=Dv_*(t+\theta,x)$,
$\bar m(t,x)=m_*(t+\theta,x)$, and
$a(t,x)=f_m(x,\bar m(t,x))$.
Uncentered linear variations $(w,\rho)$ solve
\begin{equation}\label{lin}
 \begin{cases}
 -w_t-\nu\Delta w+b\cdot Dw=a\rho,\\
 \rho_t-\nu\Delta\rho-\diver(\rho b+\bar m Dw)=0,
 \qquad\int\rho=0.
 \end{cases}
\end{equation}

\begin{proposition}\label{symplectic}
If $(w_1,\rho_1)$ and $(w_2,\rho_2)$ are two $C^{1,2}$ solutions
of \eqref{lin} on an interval, then
\[
 \Omega(t)=\int_{\T^d}(w_1\rho_2-w_2\rho_1)dx
\]
is constant. For a single solution,
\[
 \frac{d}{dt}\int_{\T^d}w\rho\,dx
 =-\int_{\T^d}a\rho^2dx-\int_{\T^d}\bar m|Dw|^2dx.
\]
If the reference solution is differentiable one additional time in the
classical sense, $(v_{*,t}(t+\theta),m_{*,t}(t+\theta))$ solves
\eqref{lin}; it is the phase variation.
\end{proposition}
\begin{proof}
Substitute \eqref{lin} in the derivative of $\int w_1\rho_2$.
The two Laplace terms cancel after spatial integration by parts, and the
terms involving $b$ cancel as well. The result is
\[
 \frac{d}{dt}\int w_1\rho_2
 =-\int a\rho_1\rho_2-\int\bar m Dw_1\cdot Dw_2.
\]
The right side is symmetric in the two solutions, proving conservation
of $\Omega$. Taking identical solutions gives the second formula.
Differentiating the reference PDEs in time, with constant $\lambda_*$,
gives the phase-variation assertion directly.
\end{proof}

The coefficient $a$ can be negative, so the displayed quadratic form is
not generally dissipative. No bounded forward solution operator or
Floquet monodromy for the full PDE is inferred from this identity.
The backward heat component of \eqref{lin} requires that distinction.

\section{A conditional nonlinear theorem with an explicit selection}
Here are the precise spaces and additional regularity assumptions.
Fix $0<\alpha<1$, assume the reference solution and $f$ are smooth,
and fix $\zeta>0$ with $\bar m\geq2\zeta$ at every phase.
For $T\geq1$ let $X_T$ be the Banach space of pairs $(w,\rho)$ of
parabolic H\"older class $C^{1+\alpha/2,2+\alpha}$ on
$[0,T]\times\T^d$, both of spatial mean zero at every time.
Use the sum of their usual parabolic norms, with H\"older seminorms
computed on time intervals of length at most one and their supremum over
such intervals. This norm controls $\|Dw\|_\infty+\|\rho\|_\infty$.
The local-in-time norm convention avoids a horizon factor in product
estimates. It is equivalent to the usual norm for each fixed $T$.

Let $Y_T$ be the product of two spatially mean-zero residual spaces
$C^{\alpha/2,\alpha}$, with the same interval convention, and two
mean-zero boundary spaces $C^{2+\alpha}(\T^d)$.
Write $\Pi r=r-\int_{\T^d}r$ for spatial mean projection and define
\begin{equation}\label{L}
 L_{T,\theta}(w,\rho)=
 \left(\Pi[-w_t-\nu\Delta w+b\cdot Dw-a\rho],\,
 \rho_t-\nu\Delta\rho-\diver(\rho b+\bar m Dw),\,
 \rho(0),\,w(T)\right).
\end{equation}
Centering the first equation matters: an uncentered constant residual
would produce a value correction proportional to $T-t$, which has
no uniform bound in a norm measuring the value itself.

\begin{lemma}[Uniform quadratic remainder]\label{remainder}
Under the smoothness and positivity assumptions just stated, there are
$r_0>0$ and $C<\infty$, independent of $T\geq1$ and $\theta$, such
that $r_0\leq\zeta$ and the nonlinear residual
\[
 N_{T,\theta}(w,\rho)=
 \left(\Pi\left[\tfrac12|Dw|^2-
   \{f(x,\bar m+\rho)-f(x,\bar m)-a\rho\}\right],
       -\diver(\rho Dw),0,0\right)
\]
satisfies, whenever $\|x\|_{X_T},\|y\|_{X_T}\leq r_0$,
\[
 \|N(x)\|_{Y_T}\leq C\|x\|_{X_T}^2,
 \qquad
 \|N(x)-N(y)\|_{Y_T}
 \leq C(\|x\|_{X_T}+\|y\|_{X_T})\|x-y\|_{X_T}.
\]
Here $x$ and $y$ denote pairs of perturbations, not spatial points.
\end{lemma}
\begin{proof}
On the indicated ball, $\bar m+\rho\geq\zeta$ and all densities lie
in a fixed compact positive interval. The derivatives of $f$ needed
below, through third order in $m$ and the corresponding spatial H\"older
bounds, are uniformly bounded on that interval. Reference norms are
uniform in phase and horizon because the reference is smooth and periodic.
Taylor's integral formula gives
\[
 f(x,\bar m+\rho)-f(x,\bar m)-a\rho
 =\rho^2\int_0^1(1-s)f_{mm}(x,\bar m+s\rho)ds.
\]
The multiplier on the right has uniformly bounded H\"older norm on
$r_0$ balls and is Lipschitz in $\rho$ in that norm, by the bounded
third derivative and the composition rule obtained by the mean-value
formula. For scalar functions the elementary product estimate
$[uv]_\alpha\leq\|u\|_\infty[v]_\alpha+
\|v\|_\infty[u]_\alpha$ applies to space and time differences.
Together with $u^2-v^2=(u-v)(u+v)$, it proves both estimates for the
Taylor remainder and $|Dw|^2$.
For the second component expand
$\diver(\rho Dw)=D\rho\cdot Dw+\rho\Delta w$ and use the same
product estimate, followed by subtraction of the two products for the
Lipschitz bound. The $X_T$ norm controls every derivative used.
Projection $\Pi$ is bounded in each residual norm. All estimates are
local on time intervals of length at most one, so taking their supremum
introduces no $T$ factor. Increasing one finite constant $C$ gives the
announced bounds.
\end{proof}

\begin{theorem}[Selected orbital shadowing from a uniform inverse]\label{conditional}
In addition to the preceding assumptions, suppose every
$L_{T,\theta}:X_T\to Y_T$ is a bounded linear isomorphism and
\begin{equation}\label{inverse}
 \sup_{T\geq1,\,\theta\in[0,P)}\|L_{T,\theta}^{-1}\|\leq K<\infty.
\end{equation}
Let $C,r_0$ be as in Lemma~\ref{remainder}. For boundary data define
\[
 d_0=m_0-m_*(\theta),\qquad
 d_T=\Pi(h-v_*(T+\theta)),\qquad
 \delta=\|d_0\|_{C^{2+\alpha}}+\|d_T\|_{C^{2+\alpha}},
\]
with $\int m_0=1$. If
\begin{equation}\label{small}
 \delta\leq\min\{r_0/(2K),\,1/(8CK^2)\},
\end{equation}
then there is exactly one centered perturbation $(w,\rho)$ in the
closed $X_T$ ball of radius $2K\delta$ producing a classical
finite-horizon MFG solution with $m(0)=m_0$, $u(T)=h$.
It satisfies
\[
 \|w\|_{X_T\text{ component}}+\|\rho\|_{X_T\text{ component}}
 \leq2K\delta,\qquad m\geq\zeta,
\]
and, for the orbit distance in the catalogue,
\[
 \sup_{0\leq t\leq T}
 \inf_{s\in[0,P)}\bigl(\|Du(t)-Dv_*(s)\|_\infty
                       +\|m(t)-m_*(s)\|_1\bigr)
 \leq2K\delta.
\]
Selection is defined by this unique small fixed point. For two boundary
perturbations satisfying a common smallness bound, their selected centered
solutions differ by at most $2K$ times the $Y_T$ norm of the difference
of the boundary vectors.
\end{theorem}
\begin{proof}
Put $d=(0,0,d_0,d_T)$ and consider the map
$\mathcal A(x)=L_{T,\theta}^{-1}(d-N_{T,\theta}(x))$.
For $R=2K\delta\leq r_0$ and $\|x\|\leq R$,
\[
 \|\mathcal A(x)\|\leq K\delta+KCR^2
 \leq K\delta+\tfrac12K\delta<R
\]
when $\delta>0$, by \eqref{small}; for $\delta=0$ the zero vector
is the sole point of the stated ball. For $\|x\|,\|y\|\leq R$,
\[
 \|\mathcal A(x)-\mathcal A(y)\|
 \leq2KCR\|x-y\|\leq\tfrac12\|x-y\|.
\]
The complete closed ball therefore has a unique fixed point by the
contraction theorem (whose proof is the convergence of iterates under
this geometric distance bound). It solves $Lx+N(x)=d$.
The density $m=\bar m+\rho$ has mass one and is at least
$2\zeta-r_0\geq\zeta$.

To recover the actual value, set
$u_{\rm ref}(t,x)=v_*(t+\theta,x)-\lambda_*t$.
The first residual of $u_{\rm ref}+w$ with density $m$ has zero
spatial projection and is therefore a scalar function $r(t)$.
Choose a scalar $c(t)$ by
$c'(t)=r(t)$ and
$c(T)=\int_{\T^d}(h-u_{\rm ref}(T))dx$.
Then $u=u_{\rm ref}+w+c$ satisfies the unprojected HJB equation
because its residual is $r-c'=0$, and its terminal value is $h$.
The Fokker--Planck equation and initial condition already follow from
the fixed-point equation. This also proves nonemptiness for every
allowed perturbation; no vacuous stability bound is used.

At time $t$, use orbit phase $s=t+\theta$ modulo $P$.
The difference of gradients is $Dw$ and the difference of densities is
$\rho$. The stated $X_T$ norm and unit torus volume bound their required
norms by $\|x\|\leq2K\delta$, proving shadowing.
For two fixed points, subtract their fixed-point equations and use the
same $1/2$ Lipschitz bound for the nonlinear part to obtain
$\|x-x'\|\leq K\|d-d'\|+\tfrac12\|x-x'\|$.
This proves the final assertion and completes the proof.
\end{proof}

This theorem measures perturbations in a stronger boundary topology than
the catalogue's $L^1$ and gradient-uniform topology and selects the nearby
branch. It neither controls far-away equilibria with the same data nor
verifies \eqref{inverse} for a genuinely periodic reference. These are
explicit restrictions, not suppressed steps.

\section{Exact singular horizons in the homogeneous linear model}
\begin{proposition}\label{resonant}
Linearize the homogeneous quadratic MFG around $v=0,m=1$, with
$a=f'(1)$. Let $q>0$ be a spatial Laplace eigenvalue and suppose
$a<-\nu^2q$. Put $\omega=\sqrt{-q(\nu^2q+a)}$.
For the corresponding mode amplitudes $(U,R)$, impose $R(0)=0$ and
$U(T)=d$. Define
\[
 D(T)=\cos(\omega T)+\frac{\nu q}{\omega}\sin(\omega T),\qquad
 T_j=\frac{j\pi-\arctan(\omega/(\nu q))}{\omega},\quad j=1,2,\ldots.
\]
If $D(T)\ne0$, the unique mode solution has $U(0)=d/D(T)$.
At every $T_j$, zero boundary data have a nonzero linear solution,
and boundary data with $d\ne0$ have no linear solution.
Furthermore $|D'(T_j)|=\sqrt{-aq}>0$, so the boundary response gain
is unbounded as $T\to T_j$. Hence a bound of the form
\eqref{inverse} over all $T\geq1$ fails for this homogeneous reference.
\end{proposition}
\begin{proof}
The mode system is
\[
 \binom{U'}{R'}=
 A\binom U R,\qquad
 A=\begin{pmatrix}\nu q&-a\\-q&-\nu q\end{pmatrix},
 \qquad A^2=-\omega^2 I.
\]
Thus $e^{tA}=I\cos(\omega t)+(A/\omega)\sin(\omega t)$.
Starting at $(U(0),0)$ gives
\[
 U(t)=D(t)U(0),\qquad R(t)=-\frac q\omega\sin(\omega t)U(0).
\]
The terminal condition therefore is the scalar equation $D(T)U(0)=d$.
The listed $T_j$ are exactly the positive zeros of $D$.
At a zero, $\cos(\omega T_j)=-(\nu q/\omega)\sin(\omega T_j)$,
and $\sin(\omega T_j)\ne0$. Direct differentiation gives
\[
 |D'(T_j)|=\sqrt{\omega^2+\nu^2q^2}=\sqrt{-aq}.
\]
This proves both the singularity and the divergent inverse gain.
There are arbitrarily large $T_j$, so restricting to $T\geq1$ does
not remove the obstruction. A smooth spatial eigenfunction embeds each
mode in all the H\"older spaces used above, so it obstructs a bounded
inverse for the PDE operator as well.
\end{proof}

\section{Exact remaining issue}
The homogeneous state in Proposition~\ref{resonant} is stationary;
that proposition is not a nonlinear instability theorem for an existing
genuine periodic orbit. It demonstrates why a horizon-independent
boundary inverse cannot be presumed from imaginary mode eigenvalues.
For the genuinely periodic solutions in the target, one must establish an
appropriate uniform inverse or another nonlinear selection mechanism,
handle phase and possible conjugate horizons, and, for the strongest
catalogue formulation, control every admissible equilibrium in the weaker
boundary topology. Those steps are not proved here. The conditional
selected shadowing theorem, the conserved form, and the exact singular
horizons are the verified outcomes of this attempt.

\begin{thebibliography}{9}
\bibitem{CD} P. Cardaliaguet and F. Delarue,
\emph{Selected topics in mean field games}, Proceedings of the ICM 2022,
vol. 5, Section 3.3. DOI \href{https://doi.org/10.4171/ICM2022/17}{10.4171/ICM2022/17}.
\url{https://ems.press/content/book-chapter-files/33265}.
\bibitem{CN} M. Cirant and L. Nurbekyan,
\emph{The variational structure and time-periodic solutions for mean-field games systems},
Minimax Theory and its Applications 3 (2018), 227--260.
\url{https://arxiv.org/pdf/1804.08943v1}. Section 6.
\end{thebibliography}
\end{document}
